\section{原生多模态 vs 桥接式多模态}

\subsection{两种路线的对比}

主持人引入讨论：原生多模态（直接将 pixel 变成 token，从早期就图文混合训练）vs 桥接式（先训 LLM，再接 ViT + adapter）。

\begin{importantbox}{薛福昭：Native 更 Clean，但 Encoder 有工程优势}
\textbf{桥接式（Encoder-based）的优点}：
\begin{itemize}
  \item Encoder 和 Decoder 可分开训练，更容易优化
  \item 可以灵活压缩 Token 数量，全局计算效率更好
\end{itemize}
\textbf{桥接式的缺点}：
\begin{itemize}
  \item Encoder 的训练目标（CLIP? Detection? Reconstruction?）无法保证对下游完全无损
  \item Infra 复杂度高：Image 位置不规则导致 Sequence Parallelism 时 workload 不均衡，需要复杂的 reshuffle/recompute，对 MFU 伤害很大
\end{itemize}
\textbf{Native 的优点}：pixel token 与 text token 无区别，Infra 层面几乎不需要修改纯 LM 的基础设施。
\end{importantbox}

\textbf{刘子伟（偏好 Native）}：Native 的三大优势：
\begin{enumerate}
  \item \textbf{Data 全局考虑}——不用先训 ViT 再接 LM，数据配比更统一
  \item \textbf{Emerging Ability}——早期就有 Interleaved Data，Early Fusion 可能带来涌现能力（Gemini 3 据报道重训了原生版本）
  \item \textbf{Infra Clean}——可直接复用 LM 的 Infra，部署更简单
\end{enumerate}
但目前缺乏一个所有人都能用的开源 Native Multimodal Codebase/Checkpoint。

\textbf{谢天宝}：Native Multimodal 从 2023 年 Gemini 1 就开始讨论。本质上是工程成熟度的问题——工业界手里有想法但工程能力还没完全跟上，还在吃当前方案的红利。\textbf{最终哪种方案留下来，取决于可维护性和可控性}（类比软件工程）。

\textbf{曹玉（数据视角）}：多模态模型的发展\textbf{跟着数据的可获得性走}。原生多模态有更好的数据亲和力——对于后训练来说，base model 基础能力越强越好，native 架构在这方面有天然优势。同时提出一个未解问题：原生多模态的\textbf{输出端}应该包含哪些模态？输入端（text/vision/audio）没争议，但输出端各执一词——人本身并没有 vision native 输出能力。

\subsection{LLM 与 VLM 为何分开？何时统一？}

主持人指出：国内很多厂商仍将 LLM 和 VLM 分开发布，而国外（GPT、Gemini）倾向统一模型。

\textbf{薛福昭}：闭源模型是黑盒，可能内部就是两个模型路由。如果用 Encoder + Decoder 结构部署，最佳做法是分两个芯片——text-only query 不需要 vision encoder 的 memory。合在一起会让问题更复杂，但在 Serving 层面是可解的。

\textbf{谢天宝}：LLM/VLM 分开与 Coding/Math/Tool Use 分开研发的原因一样——分支开发、然后合并是工程上的常见做法。目前处于一个``中间版本''的状态。

\textbf{曹玉}：VLM 和 LLM 分开是暂时的。对于真正有智能的操作体，理解世界的能力\textbf{一定不是单一模态的}。当前分离是因为 Transformer 在 discrete signal 上远比 continuous signal 建模成熟。长期一定会克服工程/组织/技术问题，迈向原生输入多模态。

\textbf{刘子伟}：短期分开有历史因素（Benchmark 偏文字智能、Vision 对排名影响不大）。但从 Physical AI 的角度看——需要长程决策（long-horizon action/decision）和多模态融合，这些场景\textbf{必须统一}。人类是``先有运动能力后有语言''，AI 是反过来，最终也要解那些需要 sensory input 的底层任务。

\subsection{本章小结}
Native vs Encoder-based 各有优劣，短期共存，长期趋向 Native。LLM/VLM 分开是工程/组织的暂时状态，Physical AI 时代必然统一。

%% ===== 合成数据与 Mid-training =====
